\chapter{Organizing larger applications}

\noindent\textbf{Learning path: Complete path / scale-up chapter.} Read when a single-module application starts becoming difficult to change safely.\par\medskip

\rustbridgeblock{Module decomposition, qualified names, explicit dependencies, and keeping related code together follow familiar Rust design instincts.}{Velran resolves an application-root source graph, adds domain \code{object} namespaces, and keeps HTTP \code{route} policy separate from source/module structure.}{Do not relearn modularity. Use familiar module boundaries; learn only the extra distinction between source location, domain namespace, and external HTTP exposure.}


As an application grows, change locality matters more than file count. Use \code{mod} for source structure and \code{object} for domain API namespaces; keep HTTP exposure in \code{route}.

\section{Three separate questions}

Treat the structure of a larger application as three distinct layers:

\begin{center}
\begin{tabularx}{\textwidth}{@{}l Y@{}}
\toprule
\textbf{Tool} & \textbf{Responsibility} \\
\midrule
\code{mod} & In which file or source module does the code live? \\
\code{object} & Which business concept owns the model, query, page, or action? \\
\code{route} & What HTTP surface and what auth/cache/rate-limit policy expose it? \\
\bottomrule
\end{tabularx}
\end{center}

This separation is deliberate. A domain object does not automatically receive a database, current user, filesystem, or network access, and a module does not imply HTTP policy.

\section{\code{mod} is not textual include}

A small multi-file project might look like:

\begin{lstlisting}
myapp/
|-- main.vrn
|-- models.vrn
|-- queries.vrn
|-- components.vrn
|-- pages.vrn
|-- actions.vrn
|-- migrations/
`-- public/
\end{lstlisting}

The entry point:

\begin{lstlisting}[language=Velran]
mod models;
mod queries;
mod components;
mod pages;
mod actions;
\end{lstlisting}

\code{mod} is not C-style textual inclusion and it does not inject declarations into a global namespace. The compiler loads, validates, and compiles the explicit source graph into one program. A module owns a namespace: declarations from \code{queries.vrn}, for example, are referenced as \code{queries::name} across module boundaries. Module paths are application-root relative, so nested source files must be loaded explicitly with paths such as \code{mod article::admin;}.

\section{When should you switch to feature-based structure?}

Splitting by technical type -- \code{models.vrn}, \code{queries.vrn}, \code{pages.vrn} -- is a good first step, but in a larger application it can mean opening five or six files for one feature change. At that point, organize by business area instead:

\begin{lstlisting}
main.vrn
article.vrn
article/
|-- admin.vrn
`-- components.vrn
account.vrn
account/
`-- profile.vrn
invoice.vrn
invoice/
`-- admin.vrn
routes.vrn
migrations/
public/
\end{lstlisting}

Then \code{main.vrn} becomes mostly application composition:

\begin{lstlisting}[language=Velran]
mod article;
mod article::admin;
mod article::components;
mod account;
mod account::profile;
mod invoice;
mod invoice::admin;
mod routes;
\end{lstlisting}

The goal is change locality: one business change should stay mostly inside one feature directory.

\section{Domain object: a business namespace, not an OOP class}

\code{object} groups related business names inside a domain namespace. For example:

\needspace{30\baselineskip}
\begin{lstlisting}[language=Velran]
object Article {
    model {
        id: i64
        slug: Slug
        title: String
        body: String
        authorUsername: String
    }

    #[query]
fn bySlug(db: Db, slug: Slug)
        -> Result<Article, DbError> sql {
        SELECT id, slug, title, body, authorUsername
        FROM articles
        WHERE slug = :slug
    }

    #[page]
fn show(ctx: PageContext, db: Db, slug: Slug)
        -> Result<Html, PageError> {
        let article = Article.bySlug(db, slug)?;
        let markdown_html = commonmark::render(&article.body);
        return Ok(html {
            <article>
                <h1>{{ article.title }}</h1>
                {{ markdown_html }}
            </article>
        });
    }
}
\end{lstlisting}

Usage becomes expressive:

\begin{lstlisting}[language=Velran]
let article = Article.bySlug(db, slug)?;
\end{lstlisting}

Short domain names become natural too, for example \path{Invoice.approve}, \path{Invoice.cancel}, and \path{UserAccount.disable}.

\subsection*{What may an object contain?}

A framework domain object can group operations around one model. The following is only a structural sketch, not standalone compilable Velran source:

\begin{lstlisting}
object Article {
    model { ... }
    #[query]
fn bySlug(...) ...
    #[page]
fn show(...) ...
    #[action]
fn publish(...) ...
    #[component]
fn Card(...) ...
    #[layout]
fn EditorLayout(...) ...
}
\end{lstlisting}

Exactly one model block belongs to a framework domain \code{object}. This construct groups framework-owned model/query/page/action names; it is separate from the Rust-like \code{struct} + inherent \code{impl} surface used for ordinary pure/domain values. Framework objects still do not imply inheritance, hidden capabilities, virtual dispatch, or reflection. For Rust-like values, associated constructors (for example \code{Type::new}) and immutable \code{&self} methods are supported; general mutable/owned receivers remain intentionally fail-closed.


\section{Package boundaries for reusable application code}

When code must be reused across applications or isolated behind a stable library boundary, move it into a local package instead of stretching the application module graph indefinitely. Each package has its own \code{velran.toml}, private-by-default API, and bounded local dependency graph. Direct dependencies are explicit path dependencies; transitive dependencies stay encapsulated rather than becoming ambient application namespaces.

Use \code{pub struct}, \code{pub fn}, public inherent methods, and the narrow package-root module re-export when designing the library surface. Keep framework authority in the application: a reusable library cannot gain HTTP routes, actions, database/network authority, or other web capabilities merely through \code{pub}. This makes package decomposition an architectural boundary, not a capability-escalation mechanism.

\section{Dependencies remain explicit}

The domain namespace does not hide capabilities. This is good:

\needspace{11\baselineskip}
\begin{lstlisting}[language=Velran]
#[query]
fn byId(db: Db, id: i64)
    -> Result<Article, DbError> sql {
    SELECT id, slug, title, body, authorUsername
    FROM articles
    WHERE id = :id
}
\end{lstlisting}

An action receives the capabilities it needs just as explicitly. The following is intentionally abbreviated structural code; the concrete typed state-changing query is shown in the CRUD and business-audit sections:

\needspace{12\baselineskip}
\begin{lstlisting}
#[action]
fn publish(ctx: ActionContext, db: Db, id: i64)
    -> Result<Redirect, PageError> {
    let article = Article.byId(db, id)?;
    authorize article owner authorUsername or role Publisher;
    // ... state transition inside a transaction ...
    return Ok(redirect(adminArticles()));
}
\end{lstlisting}

The handler signature still reveals database access. The same principle applies to AppFs, auth, and named network egress: a domain object must not become a hidden service locator.

\section{Why should routes remain top-level?}

Routes deliberately do not disappear entirely into objects:

\begin{lstlisting}[language=Velran]
route articleShow GET "/article/:slug<Slug>"
    public cache public ttl 60
    => article::Article.show;

route articlePublish POST "/admin/article/:id<i64>/publish"
    auth role Publisher
    invalidate cache articleShow
    => article::Article.publish;
\end{lstlisting}

A route registry or code review can then see in one place:

\begin{itemize}
  \item URL and HTTP method;
  \item authentication/authorization requirement;
  \item cache policy and invalidation;
  \item rate limit or resource profile;
  \item the global route name used in logs and audit.
\end{itemize}

A domain member name may be short and locally meaningful, such as \code{Article.publish} inside its own module; from a separate \code{routes.vrn}, use \code{article::Article.publish}. The route name should remain globally descriptive, such as \code{articlePublish}.

\section{Recommended naming convention}

\begin{center}
\begin{tabularx}{\textwidth}{@{}l l Y@{}}
\toprule
\textbf{Element} & \textbf{Form} & \textbf{Example} \\
\midrule
object & PascalCase business noun & \code{Article}, \code{Invoice}, \code{UserAccount} \\
member & camelCase, short inside domain & \code{bySlug}, \code{publish}, \code{approve} \\
route & globally auditable name & \code{articleShow}, \code{invoiceApprove} \\
module & feature or technical unit & \code{article}, \code{account}, \code{routes} \\
\bottomrule
\end{tabularx}
\end{center}

\section{One possible CMS structure}

A medium-size CMS might use:

\needspace{18\baselineskip}
\begin{lstlisting}
main.vrn
routes.vrn
article.vrn            # Article object + base queries
article/
|-- admin.vrn          # edit/publish
`-- components.vrn     # Card, lists
media.vrn
media/
`-- admin.vrn
account.vrn
migrations/
public/
`-- app.css
\end{lstlisting}

Public and admin HTTP policy remain in the route layer. The article feature keeps the article business model and its operations together.

\section{What not to do}

\subsection*{Do not create one file per declaration}

Overly small modules create more navigation than structure. A fifty-line feature does not need eight files.

\subsection*{Do not build a hidden ``service'' layer}

If an operation uses DB, AppFs, or auth context, keep that visible in the typed signature. One of Velran's strengths is precisely that capability boundaries are explicit.

\subsection*{Do not duplicate route policy inside handlers}

Route-level auth, cache, rate limit, and resource policy are declarative contracts. Domain-specific object authorization naturally remains in the handler because it depends on the business record.

\subsection*{Do not use free-form String for business state}

If a domain has \code{Draft}, \code{Review}, \code{Published}, or \code{Cancelled} states, use an enum. Project structure helps most when the domain contract is typed too.

\section{Refactoring path: from one file to a larger application}

You do not need to design the final directory hierarchy on day one. A good evolutionary path is:

\begin{enumerate}
  \item Start with one \code{main.vrn} while the application is small.
  \item Separate model/query/page/action groups when the file becomes hard to navigate.
  \item If one feature is scattered across many technical files, switch to feature-based modules.
  \item If too many global business names appear, group them under \code{object} namespaces.
\end{enumerate}

\section{Verified companion examples}
For module structure, use \path{examples/module-namespaces/main.vrn}; for a generic multi-file starting point, use \path{examples/starter-project/main.vrn}; and for the running project in this book, use \path{examples/secure-notes/main.vrn}. All three are compiler-checked and therefore better copy/paste sources than shortened architectural sketches.

\section{Practice: refactor by change locality}
\paragraph{Exercise mode: independent.} Use the independent method defined in the preface.

Pick one feature and identify every file that changes when the feature changes. If a small business change requires edits across many unrelated technical layers, consider regrouping the code around the feature or domain object. Refactor only after the dependency direction is clear.

\section{Common mistake: modularizing by file count alone}
Useful modules own coherent contracts and reduce concepts that must be held at once; avoid both giant files and artificial one-function modules.

\section{Running project checkpoint: Secure Notes}
Move Secure Notes from the initial single-file shape into a structure that keeps note-specific models, queries, pages/actions, and policies easy to find. Preserve the same externally visible routes and security behavior while improving change locality.
